\textit{In the setting of Section 1.46, prove that 
  %
    $f \mapsto \deriv[\alpha]{f}$
  %
is a continuous mapping of 
%
  $C^{\infty}\left(\Omega\right)$ into 
  $C^{\infty}\left(\Omega\right)$ and also of 
  $\testFunctions{K}$ into 
  $\testFunctions{K}$, for every multi-index $\alpha$.
%
}
\begin{proof} 
In both cases, $\deriv[\alpha]$ is a linear mapping.
It is then sufficient to establish continuity at the origin.
%
We begin with the $C^\infty\left(\Omega\right)$ case.
\skipline
Let $U$ be an arbitrary neighborhood of the origin.
It exists $N$ such that $U$ contains
%
  \begin{equation} 
    V_{N}= \left\{
      \phi \in C^\infty\left(\Omega\right): 
      \max\set{
        | \deriv[\beta]{\phi}(x) |
      }{
        \magnitude*{ \beta } \leq N, x\in K_N
      }
    < 1/N
    \right\}.
  \end{equation}
%
Now pick $g$ in $V_{N+\magnitude*{\alpha}}$so that
%
  \begin{equation}
    \max
    \set{
      \magnitude*{ \deriv[\gamma]{g}\left(x\right) }
    }{
      \magnitude*{ \gamma | \leq N+| \alpha }, 
      x\in K_{N}
    }
    < \frac{1}{N+\magnitude*{\alpha}}.
  \end{equation}
%
(the fact that $K_N\subset K_{N+\magnitude*{\alpha}}$ was tacitly used).
%
The special case $\gamma = \beta + \alpha$ yields
%‡
\begin{equation}
    \max
    \set{
      | \deriv[\beta]{\deriv[\alpha]{g}}(x)|
    }{
      \magnitude*{\beta} \leq N, 
      x\in K_{N}
    }
    < \frac{1}{N}.
  \end{equation}
%
We have just proved that
%
  \begin{equation}\label{1.17. inclusion}
    g \in V_{N + \magnitude*{\alpha}}
      \then
    \deriv[\alpha]{g} \in V_{N},
      \qquad
      \text{, \ie, }
      \qquad
    \deriv[\alpha] \left(V_{N+\magnitude*{\alpha}}\right) \subset V_N,
  \end{equation}
%
which establishes the continuity of %
%
  $\deriv[\alpha]: C^\infty \left(\Omega\right)\to C^\infty \left(\Omega\right)$. %
%
\skipline
To prove the continuity of the restriction 
%
  $\deriv[\alpha] \lvert_{\testFunctions{K}}: \testFunctions{K} \to \testFunctions{K}$, %
%
%
we first remark that the collection of the  
%
  ${V_N \cap \testFunctions{K}}$ % 
%
is a local base of the subspace topology of $\testFunctions{K}$.
%
%
  $V_{N+\magnitude*{\alpha}} \cap \testFunctions{K}$ % 
%
is then a neighborhood of $0$ in this topology. %
Furthermore, 
%
  \begin{align}
    %
    \deriv[\alpha] \lvert_{\testFunctions{K}} \left(V_{N+\magnitude*{\alpha}} \cap \testFunctions{K}\right)
    %
    & =
      \deriv[\alpha]\left(V_{N+\magnitude*{\alpha}} \cap \testFunctions{K}\right) \\
      %
    & \subset
      \deriv[\alpha]\left(V_{N+\magnitude*{\alpha}}\right)
        \cap
      \deriv[\alpha]\left(\testFunctions{K}\right) \\
      %
    & \subset 
      V_N 
        \cap 
      \testFunctions{K}
        %
          \qquad (\text{see }(\ref{1.17. inclusion}))
        %
    %
  \end{align}
%
So ends the proof.
\end{proof}
